\documentclass[12pt,a4paper]{article}
\usepackage[margin=2.5cm]{geometry}
\usepackage[utf8]{inputenc}
\usepackage[T1]{fontenc}
\usepackage[russian]{babel}
\usepackage{hyperref}
\usepackage{graphicx}
\usepackage{booktabs}
\usepackage{subcaption}
\usepackage{caption}
\usepackage{siunitx}

\title{Промежуточный отчёт: ускорение инференса Super-Resolution U-Net}
\author{Команда 52: Артём Чубов, Андрей Петухов, Евгений Веселков}
\date{1 апреля 2026}

\begin{document}

\maketitle

\section{Введение}
Цель проекта состоит в обучении и последующем ускорении инференса U-Net-подобной модели. В качестве задачи, которую должна решать наша модель, мы выбрали single-image super-resolution, мы сравниваем несколько вариантов инференса на одной и той же архитектуре и одном и том же наборе данных. В работе оцениваются как метрики качества восстановления изображения, так и вычислительные характеристики. Основное внимание уделяется метрикам \href{https://en.wikipedia.org/wiki/Peak_signal-to-noise_ratio}{PSNR} (Peak signal to noise ratio) и \href{https://en.wikipedia.org/wiki/Structural_similarity_index_measure}{SSIM} (Structural similarity index measure) с точки зрения качества, а с точки зрения инференса - latency и пропускная способность.

\section{Постановка задачи}
В проекте рассматривается задача single-image super-resolution. В качестве базовой модели используется лёгкая U-Net-подобная encoder-decoder архитектура со skip-connection-ами. Эксперименты проводятся на датасете DIV2K, разделённом на обучающую, валидационную и тестовую выборки. К настоящему моменту мы успели поэкспериментировать с коэффициентами масштабирования $\times 4$ и $\times 8$. Вход в нейросеть формируется из изображения в высоком разрешении через bicubic downsampling, после чего он bicubic upsample-ится обратно до исходного размера и в таком виде подаётся в модель. Качество реконструкции оценивается с помощью PSNR и SSIM. Эксперименты организованы в виде единого репозитория проекта на \href{https://github.com/aapetukhov/efml-unet}{GitHub}.

\section{Ход первого этапа}
На первом этапе была подготовлена схема работы с данными, реализован и обучен бейзлайн, после чего проведена первичная оценка качества на валидации. Следующим шагом стали замеры производительности инференса на GPU в режимах FP32 и FP16. На данный момент сравнения FP32 и FP16 проведены для модели $\times 4$. Для масштаба $\times 8$ получены первые результаты по качеству на валидации.

\section{Текущие результаты}

Промежуточные результаты первого эксперимента для модели U-Net SR со scale factor $\times 4$ приведены в таблице~\ref{tab:fp32_fp16_comparison}.

\begin{table}[h!]
\centering
\caption{Сравнение FP32 и FP16 инференса для U-Net SR $\times 4$}
\begin{tabular}{lcccccc}
\toprule
Пайплайн & Latency, ms & P50, ms & P95, ms & Throughput, img/s & PSNR & SSIM \\
\midrule
FP32 & 5.71 & 5.35 & 6.70 & 175.2 & 23.86 & 0.678 \\
FP16 & \textbf{5.15} & \textbf{5.02} & \textbf{6.23} & \textbf{194.2} & 23.86 & 0.678 \\
\bottomrule
\end{tabular}
\label{tab:fp32_fp16_comparison}
\end{table}

\begin{figure}[h]
    \centering
    \begin{subfigure}[t]{0.48\textwidth}
        \centering
        \includegraphics[height=5cm]{pics/0802_lr.jpeg}
        \caption{Входное изображение в низком разрешении}
    \end{subfigure}\hfill
    \begin{subfigure}[t]{0.48\textwidth}
        \centering
        \includegraphics[height=5cm]{pics/0802_prediction.png}
        \caption{Предсказание модели, разрешение $\times 8$}
    \end{subfigure}
    \caption{Визуальный пример работы бейзлайн-модели}
    \label{fig:sr_example}
\end{figure}

На рисунке~\ref{fig:sr_example} приведён пример работы модели для масштаба $\times 8$. Слева показан входной кадр, а справа - предсказание модели.

\section{Дальнейший план}
Следующий этап проекта связан с расширением числа сравниваемых режимов инференса. После фиксации бейзлайна в FP16 планируется исследовать ускорение с использованием \texttt{torch.compile} и сравнить несколько режимов компиляции. Затем предполагается перейти к альтернативным компиляторным стекам, в частности TVM или близким решениям, чтобы проверить, можно ли получить дополнительный выигрыш за счёт autotuning и более агрессивной оптимизации вычислительного графа.

Отдельное направление связано с исследованием квантизации, прунинга и спарсности. Здесь интересно, конечно, насколько сильно такие преобразования повлияют не только на скорость, но и на итоговое качество.


\end{document}